\documentclass[11pt]{scrartcl}
\usepackage{aura}

\title{USA TST 2026/1}


\begin{document}
\maketitle

\begin{problem*}{}{}
    Let $n$ be a positive integer. Prove that one can paint the non-zero coefficients of the polynomial
\[ f(x_1, x_2, \dots, x_n) = \prod_{k=0}^n (x_1 + x_2 + \dotsb + x_n - k) \]with $2^n - 1$ colors such that the coefficients of each color have sum $0$, and each color is used at least once.
\end{problem*}
\begin{proof}
    For ease, define the polynomials \[P(x) = x(x-1)(x-2)\dots(x-n)\]and \(Q(x_1,x_2,\dots,x_k)=P(x_1+x_2+\dots+x_k)\)
    Consider any subset \(S \subset \{1, 2, \dots, n\}\). For every such \(S\), we define the color \(C_S.\) Suppose \(S = \{s_1, s_2, \dots, s_t.\}\) Then for every term of the form \(x_{s_1}^{e_1}x_{s_2}^{e_2}\dots x_{s_t}^{e_t}\) we color its coefficient with the color \(C_S.\) Here, of course, all the \(e_i\)'s are greater than \(0\). Now clearly \(2^n - 1\) colors are used. We prove that this works.

\begin{claim*}{}{}
    The sum of coefficients of the terms colored with the color \(C_S\) is \(0\).
\end{claim*}
\begin{proof} 
    Let \(S = \{x_1, x_2, \dots, x_t\}\). The main idea is that the desired coefficients are a subset of the coefficients in the polynomial\[P(x_1 + x_2 + \dots + x_t)\]We now use strong induction on \(|S|\). For the base case, WLOG let \(S = \{x_1\}.\) Then the coefficients of \(x_1^e\) are just the coefficients in \(P(x_1)\), which clearly add up to \(0\) since \(P(1) = 0.\) Now for the inductive step, assume it holds true for all \(|S| \le t-1.\) Then WLOG let \(S = \{x_1, x_2, \dots, x_t\}\). Clearly if we consider the coefficients of the polynomial \[P(x_1 + x_2 +\dots + x_t)\] then the coefficients we are considering form a subset of these. However the remaining coefficients form all possible coefficients of the terms where \(|S| \le t-1\) and hence the sum of the coefficients under consideration is just \[P(t) - c_1P(t-1) - c_2P(t-2) \dots -c_{t-1}P(1)=0\]where all the \(c_i\)'s are positive integers. Hence done. 
\end{proof}
In fact we have \[c_i = \binom ti\]but that doesn't matter.
\end{proof}
\end{document}